\chapter{Elastic net factor regression: derivations}\label{ch:en}

This note documents the regression module of \code{nabia}: the elastic net solver
(\code{ElasticNet}) and the \code{RegressionEngine}. The engine regresses a target series on a set of candidate factors.
It chooses how many factors to use and which ones, the penalty $\lambda$ and the mix
$\alpha$, and measures the uncertainty of the result.

We first derive the elastic net from its optimality conditions:
\begin{itemize}
  \item the coordinate descent update and soft thresholding;
  \item the smallest penalty that zeroes every coefficient, $\lambda_{\max}$;
  \item the role of the intercept and of standardization;
  \item the grouping effect that makes the elastic net stable with correlated factors;
  \item the effective degrees of freedom.
\end{itemize}
We then describe the engine:
\begin{itemize}
  \item the chronological train/test split;
  \item time-series cross-validation (blocked $K$-fold and walk-forward, with a gap);
  \item the one-standard-error rule;
  \item the in-sample criteria (BIC, AIC) used when there is no test set;
  \item exhaustive and forward subset search;
  \item the moving block bootstrap.
\end{itemize}
The last sections map each step to the code, show how to use it from C++ and from Python,
and list the tests and the limitations.




%==============================================================================
\section{Problem and notation}

We observe a target $y_t$ (for example a fund's daily return) and $N$ candidate factors
$x_{t1},\dots,x_{tN}$ (market, rate, commodity returns) for $t=1,\dots,T$, in time order.
We look for a sparse linear model
\[
  y_t = b_0 + \sum_{j\in S} b_j x_{tj} + \varepsilon_t,\qquad S\subseteq\{1,\dots,N\},\quad
  k_{\min}\le|S|\le k_{\max},
\]
that \emph{extrapolates}, that is, predicts $y$ well on dates after those used to fit it.
When only in-sample explanation is wanted, the same machinery chooses the model by an
information criterion instead.

Throughout, $X\in\R^{n\times p}$ is the design of the rows and columns at hand (a training
window, a fold, a subset), $x_j$ its $j$-th column, $\norm{\cdot}$ the Euclidean norm and
$\sth(z,t)=\operatorname{sign}(z)\max(|z|-t,0)$ the soft-thresholding operator.

%==============================================================================
\section{The elastic net}

\subsection{Objective}
For $\lambda\ge0$ and $\alpha\in[0,1]$, the elastic net in the glmnet parametrization
solves
\begin{equation}\label{en:eq:enet}
  \min_{b_0,\,b}\;
  \underbrace{\frac{1}{2n}\norm{y-b_0\mathbf 1-Xb}^2}_{\text{fit}}
  +\lambda\Big(\alpha\norm{b}_1+\frac{1-\alpha}{2}\norm{b}_2^2\Big).
\end{equation}
\begin{itemize}
  \item $\alpha=1$ is the \emph{lasso}: the $\ell_1$ penalty sets coefficients exactly to
        zero, so it selects factors.
  \item $\alpha=0$ is \emph{ridge} regression: the $\ell_2$ penalty shrinks coefficients
        but never zeroes them.
  \item $\lambda=0$ is ordinary least squares.
\end{itemize}
We write $\ell_1=\lambda\alpha$ and $\ell_2=\lambda(1-\alpha)$ for the two penalty weights.
The intercept is never penalized.

\subsection{Intercept and standardization}\label{en:sec:standardize}

\begin{proposition}[Intercept]\label{en:prop:intercept}
For any $b$, the optimal intercept in \eqref{en:eq:enet} is $b_0=\bar y-\bar x^\T b$, where
$\bar y$ and $\bar x$ are the means of $y$ and of the columns of $X$. Substituting it
leaves problem \eqref{en:eq:enet} without intercept on the centred data
$\tilde y=y-\bar y\mathbf 1$, $\tilde X=X-\mathbf 1\bar x^\T$.
\end{proposition}
\begin{proof}
The objective is a convex quadratic in $b_0$ with derivative
$-\frac1n\mathbf 1^\T(y-b_0\mathbf 1-Xb)=-(\bar y-b_0-\bar x^\T b)$, which vanishes at the stated
value. Then $y-b_0\mathbf 1-Xb=\tilde y-\tilde Xb$.
\end{proof}

\paragraph{Standardization.} Write $s_j^2=\frac1n\norm{\tilde x_j}^2$ (the population variance
of factor $j$). With \code{standardize}, the solver works on $z_j=\tilde x_j/s_j$ and finds
coefficients $\beta_j$ for those columns. The coefficients in original units are
\[
  b_j=\beta_j/s_j,\qquad b_0=\bar y-\bar x^\T b.
\]
Equivalently, the penalty in original units is
$\lambda\big(\alpha\sum_js_j|b_j|+\tfrac{1-\alpha}2\sum_js_j^2b_j^2\big)$: each factor is
charged in proportion to its volatility. A factor's units, or its leverage (a 2x version of
an index), therefore do not change whether it is selected. Without \code{standardize}, the
penalty applies to $b$ directly.

\paragraph{Constant factors.} A column whose centred mean square is below
$10^{-24}\max(1,\frac1n\norm{x_j}^2)$ is constant up to rounding. It is set to zero, and its
coefficient is fixed at $0$. The relative threshold is needed because centring a constant
column leaves rounding noise of the size of its magnitude.

In the working units, let $v_j=\frac1n\norm{z_j}^2$. Then $v_j=1$ with standardization,
$v_j=s_j^2$ without it, and $v_j=0$ for a constant column. The problem the solver works on is
\begin{equation}\label{en:eq:working}
  \min_\beta\;F(\beta)=\frac1{2n}\norm{\tilde y-Z\beta}^2+\ell_1\norm\beta_1+\frac{\ell_2}2\norm\beta_2^2 .
\end{equation}

\subsection{Optimality conditions}

$F$ is convex, as the sum of a convex quadratic and two norms. With $\ell_2>0$ it is
strictly convex, so the minimizer is unique. A point $\beta$ minimizes $F$ if and only if
$0\in\partial F(\beta)$, that is, with residual $r=\tilde y-Z\beta$:
\begin{equation}\label{en:eq:kkt}
  \frac1n z_j^\T r-\ell_2\beta_j
  \begin{cases}
    =\ell_1\operatorname{sign}(\beta_j) & \beta_j\ne0,\\
    \in[-\ell_1,\ell_1] & \beta_j=0.
  \end{cases}
\end{equation}
An active factor has its correlation with the residual pinned at $\pm\ell_1$, after the
ridge shrinkage. An inactive factor has a correlation with the residual below the
threshold. The tests check \eqref{en:eq:kkt} directly on solutions.

\subsection{The coordinate update}

\begin{proposition}[Soft thresholding]\label{en:prop:cd}
Fix every coefficient but $\beta_j$, with $v_j>0$, and let $r$ be the current residual. The
minimizer of $F$ over $\beta_j$ is
\[
  \beta_j^{\text{new}}=\frac{\sth(\rho_j,\ell_1)}{v_j+\ell_2},\qquad
  \rho_j=\frac1n z_j^\T r+v_j\beta_j^{\text{old}}.
\]
\end{proposition}
\begin{proof}
Let $r^{(j)}=r+z_j\beta_j^{\text{old}}$ be the partial residual, the residual without factor
$j$. As a function of $\beta_j$, up to constants,
\[
  F=\frac1{2n}\norm{r^{(j)}-z_j\beta_j}^2+\ell_1|\beta_j|+\frac{\ell_2}2\beta_j^2
   =\frac{v_j+\ell_2}{2}\beta_j^2-\rho_j\beta_j+\ell_1|\beta_j|+\text{const},
\]
because $\frac1nz_j^\T r^{(j)}=\frac1nz_j^\T r+v_j\beta_j^{\text{old}}=\rho_j$. This is a scalar problem
$\min_u\frac a2u^2-\rho u+\ell_1|u|$ with $a>0$. Its subgradient condition is
$au-\rho+\ell_1\partial|u|\ni0$:
\begin{itemize}
  \item if $u>0$, then $u=(\rho-\ell_1)/a$, which needs $\rho>\ell_1$;
  \item if $u<0$, then $u=(\rho+\ell_1)/a$, which needs $\rho<-\ell_1$;
  \item otherwise $u=0$ is optimal, because $|\rho|\le\ell_1$ means $0\in-\rho+[-\ell_1,\ell_1]$.
\end{itemize}
In every case $u=\sth(\rho,\ell_1)/a$.
\end{proof}

After the update the residual is corrected in $O(n)$: $r\leftarrow r-z_j(\beta_j^{\text{new}}-\beta_j^{\text{old}})$.
A sweep over all factors costs $O(np)$, and no matrix is ever formed or inverted.

\paragraph{Convergence.} The non-smooth part of $F$ is separable across coordinates, and the
smooth part is a convex quadratic. Cyclic coordinate descent therefore converges to a
minimizer (Tseng, 2001). Each update cannot increase $F$, because it minimizes $F$
exactly along one coordinate.

\paragraph{Stopping rule.} A change $\Delta_j$ of coordinate $j$ moves the fitted values by
$z_j\Delta_j$, whose mean square is $v_j\Delta_j^2$. The solver stops when, over a full
sweep,
\[
  \max_jv_j\Delta_j^2<\tau\cdot\tfrac1n\norm{\tilde y}^2,
\]
that is, when no coordinate moves the fit by more than a fraction $\tau$ (\code{tolerance},
default $10^{-7}$) of the target's variance. The rule is the same in standardized and in
raw units, and does not depend on the scale of $y$.

\paragraph{Active-set iterations.} After a full sweep, the solver sweeps only the non-zero
coefficients until they settle. Then it checks with a new full sweep. It stops when a
\emph{full} sweep meets the rule, so a factor that should enter is never missed. Inactive
factors usually stay at zero, so most sweeps cost $O(n|A|)$ instead of $O(np)$.

\subsection{The smallest penalty that zeroes everything}

\begin{proposition}[$\lambda_{\max}$]\label{en:prop:lmax}
For $\alpha>0$, $\beta=0$ solves \eqref{en:eq:working} if and only if
\[
  \lambda\ge\lambda_{\max}=\frac{\max_j|z_j^\T\tilde y|/n}{\alpha}.
\]
\end{proposition}
\begin{proof}
At $\beta=0$, $r=\tilde y$ and every coefficient is inactive. By \eqref{en:eq:kkt}, $0$ is optimal
if and only if $|\frac1nz_j^\T\tilde y|\le\ell_1=\lambda\alpha$ for every $j$.
\end{proof}
For ridge ($\alpha=0$) no finite $\lambda$ zeroes the coefficients. As in glmnet, $\alpha$ is
floored at $10^{-3}$ in this formula, so ridge also gets a finite path that starts strongly
shrunk.

The code computes $\lambda_{\max}$ with the same arithmetic as the first coordinate update
and multiplies it by $1+10^{-12}$. Without that factor, rounding in the division can leave
$\lambda\alpha$ a few ulps below $|\rho_j|$, which leaves a coefficient of order $10^{-17}$
instead of an exact zero.

\subsection{Lambda grid and warm starts}
The automatic grid has $L$ (\code{n\_lambdas}) values, log-spaced from $\lambda_{\max}$ down to
$r\lambda_{\max}$ (\code{lambda\_min\_ratio}):
\[
  \lambda_i=\lambda_{\max}\,r^{\,i/(L-1)},\qquad i=0,\dots,L-1 .
\]
The path is solved in this decreasing order, and each solution starts from the previous one
(a \emph{warm start}). The first solution is $\beta=0$ exactly. Each later solution is close
to the previous one, so it takes a few sweeps. A user-supplied grid is sorted in decreasing
order for the same reason.

\subsection{Grouping effect: why the elastic net for correlated factors}

The lasso picks one factor out of a group of highly correlated ones, somewhat arbitrarily,
and may switch between them from sample to sample. The ridge term stops this.

\begin{proposition}[Grouping]\label{en:prop:grouping}
Let the factors be standardized ($\frac1n\norm{z_j}^2=1$), let $\ell_2>0$, and let factors $i,j$
have correlation $\rho_{ij}$ and coefficients of the same sign. Then
\[
  |\beta_i-\beta_j|\le\frac{\norm{\tilde y}}{\sqrt n\,\ell_2}\sqrt{2(1-\rho_{ij})}.
\]
\end{proposition}
\begin{proof}
Both coefficients satisfy \eqref{en:eq:kkt} with the same sign $s$:
$\frac1nz_i^\T r-\ell_2\beta_i=\ell_1s=\frac1nz_j^\T r-\ell_2\beta_j$. Hence
$\ell_2(\beta_i-\beta_j)=\frac1n(z_i-z_j)^\T r$, and by Cauchy--Schwarz
$|\beta_i-\beta_j|\le\frac1{n\ell_2}\norm{z_i-z_j}\norm r$. Here
$\norm{z_i-z_j}^2=2n(1-\rho_{ij})$. Also $F(\hat\beta)\le F(0)$ gives
$\norm r\le\norm{\tilde y}$.
\end{proof}
As $\rho_{ij}\to1$, the two coefficients are forced together, and they share the loading
instead of one taking it all. In the factor replication example (S\&P 500 against Nasdaq
100, correlation $0.93$), $\alpha<1$ keeps the split stable across bootstrap samples.

\subsection{Effective degrees of freedom}\label{en:sec:df}

The information criteria of \S\ref{en:sec:insample} need the number of parameters an elastic net
fit really uses. Penalized coefficients count for less than one.

\begin{proposition}[Degrees of freedom]\label{en:prop:df}
Let $A$ be the active set of the solution and $Z_A$ its columns. For a $y$ at which the
active set and the signs are locally constant,
\[
  \mathrm{df}=\mathbb 1[\text{intercept}]+\tr\!\big(Z_A(Z_A^\T Z_A+n\ell_2I)^{-1}Z_A^\T\big)
  =\mathbb 1[\text{intercept}]+\sum_i\frac{d_i}{d_i+n\ell_2},
\]
where the $d_i$ are the eigenvalues of $Z_A^\T Z_A$.
\end{proposition}
\begin{proof}
For a fit $\hat y=g(y)$, Stein's lemma gives the unbiased degrees of freedom
$\mathrm{df}=\sum_t\partial\hat y_t/\partial y_t=\tr(\partial g/\partial y)$ (Efron, 2004). On the active set, \eqref{en:eq:kkt} with fixed signs $s_A$ reads
$\frac1nZ_A^\T(\tilde y-Z_A\beta_A)-\ell_2\beta_A=\ell_1s_A$, so
\[
  \beta_A=(Z_A^\T Z_A+n\ell_2I)^{-1}(Z_A^\T\tilde y-n\ell_1s_A).
\]
The second term does not depend on $y$ locally. Hence
$\partial(Z_A\beta_A)/\partial\tilde y=H=Z_A(Z_A^\T Z_A+n\ell_2I)^{-1}Z_A^\T$. With an intercept,
$\hat y=\bar y\mathbf 1+Z_A\beta_A$ and $\tilde y=(I-\frac1n\mathbf 1\mathbf 1^\T)y$. The centred columns are
orthogonal to $\mathbf 1$, so the trace is $1+\tr H$. Diagonalizing $Z_A^\T Z_A=UDU^\T$ gives
$\tr H=\tr\big(D(D+n\ell_2I)^{-1}\big)$.
\end{proof}
Special cases:
\begin{itemize}
  \item the lasso ($\ell_2=0$) with $Z_A$ of full rank gives $\mathrm{df}=|A|$ (+1), the
        result of Zou, Hastie and Tibshirani (2007);
  \item ridge gives the familiar $\sum d_i/(d_i+n\ell_2)$;
  \item OLS gives the number of columns.
\end{itemize}
Eigenvalues below $10^{-10}\max_id_i$ are treated as zero, which handles a rank-deficient
active set.

%==============================================================================
\section{The regression engine}

\subsection{Overview}
Given $y\in\R^T$, $X\in\R^{T\times N}$ and the parameters, the engine runs these steps:
\begin{enumerate}
  \item \textbf{Split} the rows in time order into training and test sets (\S\ref{en:sec:split}).
  \item For every candidate subset $S$ (\S\ref{en:sec:search}):
  \begin{enumerate}
    \item choose $(\alpha,\lambda)$ by time-series cross-validation on the training rows
          (\S\ref{en:sec:cv});
    \item fit the elastic net on all training rows at that $(\alpha,\lambda)$;
    \item score it by the test MSE, or by an in-sample criterion when there is no test set
          (\S\ref{en:sec:scores}).
  \end{enumerate}
  \item \textbf{Keep} the subset with the lowest score. Optionally refit it on every row
        (\S\ref{en:sec:refit}) and bootstrap it (\S\ref{en:sec:bootstrap}).
\end{enumerate}

\subsection{Chronological split}\label{en:sec:split}
The training set is the first
\[
  n=\min\big(T,\lfloor fT+10^{-9}\rfloor\big)
\]
rows, with $f$ = \code{train\_fraction}. The test set is the remaining $T-n$ rows. The
$10^{-9}$ makes $0.7\times10$ count as 7 despite rounding. Rows are never shuffled.
\begin{itemize}
  \item The test set lies after every training row, so its error measures
        \emph{extrapolation}: the model faces a period it has never seen, with its own
        regimes and correlations.
  \item With $f=1$ there is no test set, and the engine works \emph{in sample}: subsets are
        compared by \code{in\_sample\_criterion}.
\end{itemize}

\subsection{\texorpdfstring{Choosing $\lambda$ and $\alpha$}{Choosing lambda and alpha}: time-series cross-validation}\label{en:sec:cv}

\paragraph{Grid.} For each $\alpha_a$ of \code{alphas}, the engine builds a $\lambda$ grid
$\lambda_{a,0}>\dots>\lambda_{a,L-1}$. It is the automatic grid from $\lambda_{\max}$ of the
subset's training rows, or the user's \code{lambdas}. The grid differs by $\alpha$, because
$\lambda_{\max}\propto1/\alpha$, so cells are compared by position. If the grid has a single
point, there is nothing to choose, and cross-validation is skipped.

\paragraph{Folds.} The $n$ training rows are cut into contiguous blocks with boundaries
$\lfloor kn/B\rfloor$. With gap $g$ (\code{validation\_gap}):
\begin{description}
  \item[Blocked $K$-fold] ($B=K$ blocks). Fold $k$ validates on block $k$ and trains on every
        row before the block minus $g$, and every row after it plus $g$.
  \item[Walk-forward] ($B=K+1$ blocks). Fold $k=1,\dots,K$ validates on block $k$ and trains on
        the rows before it minus $g$. Training always precedes validation, as in live use.
        The cost is smaller training sets in the first folds.
\end{description}
The gap removes the rows next to the validation block from training. Returns are serially
dependent, through overlapping horizons or volatility clustering, so a training row adjacent
to a validation row leaks information about it. A fold with no validation rows, or with fewer
than 2 training rows, is an error.

\paragraph{Estimate and standard error.} For every fold $k$ and grid cell $(a,i)$, the path for
$\alpha_a$ is fitted on the fold's training rows. Its validation MSE is $m_k(a,i)$. With $K$
folds:
\[
  \mathrm{CV}(a,i)=\frac1K\sum_km_k(a,i),\qquad
  \mathrm{SE}(a,i)=\sqrt{\frac{1}{K}\cdot\frac{1}{K-1}\sum_k\big(m_k(a,i)-\mathrm{CV}(a,i)\big)^2}.
\]
These are the mean of the fold errors and its standard error, treating the folds as
independent replicates. The chosen cell is $(a^\ast,i^\ast)=\argmin\mathrm{CV}$, with ties going
to the first cell, so the result is deterministic.

\paragraph{One-standard-error rule.} With \code{one\_standard\_error\_rule}, the engine stays
on row $a^\ast$ and takes the \emph{largest} $\lambda$ whose error is within one standard error
of the minimum:
\[
  i^{\ast\ast}=\min\{\,i\le i^\ast:\ \mathrm{CV}(a^\ast,i)\le\mathrm{CV}(a^\ast,i^\ast)+\mathrm{SE}(a^\ast,i^\ast)\,\}.
\]
The CV minimum is a noisy estimate. A model that is statistically as good but more strongly
penalized (sparser, more stable) is preferred (Hastie, Tibshirani and Friedman, 2009, \S7.10).

\paragraph{Final fit of the subset.} The subset is refitted on all $n$ training rows along the
grid of $\alpha_{a^\ast}$, from $\lambda_{a^\ast,0}$ down to the chosen value. This warm-starts it
exactly as during validation.

\subsection{Scores}\label{en:sec:scores}

\paragraph{Out of sample} ($T>n$). The score is the test MSE,
$\frac1{T-n}\sum_{t>n}(y_t-\hat y_t)^2$. The engine also reports the test $R^2$:
\[
  R^2_{\text{test}}=1-\frac{\sum_{t>n}(y_t-\hat y_t)^2}{\sum_{t>n}(y_t-\bar y_{\text{train}})^2}.
\]
It is measured against the \emph{training} mean, the forecast available without the model
(Campbell and Thompson, 2008). It is negative when the model does worse than that mean.

\paragraph{In sample} ($T=n$). \label{en:sec:insample}
The score is an information criterion on the training rows. Under Gaussian errors, the
log-likelihood maximized over the noise variance is
$-2\log\hat L=n\log(2\pi\,\mathrm{MSE})+n$. Dropping the constants, which are the same for
every model:
\[
  \mathrm{BIC}=n\log\mathrm{MSE}+\mathrm{df}\log n,\qquad
  \mathrm{AIC}=n\log\mathrm{MSE}+2\,\mathrm{df},
\]
with the effective df of Proposition~\ref{en:prop:df}.
\begin{itemize}
  \item BIC is consistent: with a true sparse model it selects it with probability
        $\to1$.
  \item AIC targets prediction error and admits more factors.
  \item \code{MSE} alone never pays for a factor. Among nested least-squares fits it always
        prefers the largest subset. It is there for comparison.
\end{itemize}

\subsection{Subset search}\label{en:sec:search}

With $k_{\min}$ = \code{min\_factors} and $k_{\max}$ = \code{max\_factors} ($0$ meaning $N$):
\begin{description}
  \item[Exhaustive] evaluates every subset of each size, in lexicographic order, which is
        $\sum_{k=k_{\min}}^{k_{\max}}\binom Nk$ subsets. For each size it keeps the subset with
        the lowest score. For example, $N=14$ and sizes 1--4 give $14+91+364+1001=1470$
        subsets.
  \item[Forward] is greedy. Starting from the empty set, step $k=1,\dots,k_{\max}$ adds the
        remaining factor whose addition gives the lowest score. That costs
        $\sum_{k=1}^{k_{\max}}(N-k+1)$ evaluations, which is linear in $N$ for a fixed
        $k_{\max}$. The best subset of size $k$ is recorded for $k\ge k_{\min}$. Forward
        search can miss the best subset when factors only help together, but it reaches it
        in the common case of nested good models.
  \item[Auto] uses Exhaustive if the count above is at most \code{max\_exhaustive\_subsets}
        (default 10\,000), and Forward otherwise. The count is computed in floating point and
        saturates, so it never overflows.
\end{description}
The selected model is the best of the per-size winners (\code{best\_by\_size}), by score. Every
evaluation is independent, so the subsets of one size (exhaustive) or one step (forward) are
evaluated in parallel (\S\ref{en:sec:parallel}).

\begin{remark}[Selection on the test set]
Choosing among many subsets by their test error makes the winner's test error an
\emph{optimistic} estimate of its error on new data, because the minimum of many noisy
errors is biased down. A factor can win by noise. The bootstrap of \S\ref{en:sec:bootstrap}
helps detect this: a spurious factor's interval contains zero. For a clean final number,
keep a later period out of the engine entirely.
\end{remark}

\subsection{Refit on the full sample}\label{en:sec:refit}
With \code{refit\_on\_full\_sample} and a test set, the chosen factors are refitted on all $T$
rows at the chosen $(\alpha,\lambda)$. The reported coefficients and intercept are then the
refit, which uses the most recent data. The predictions, residuals and scores remain those of
the training-rows model, so the test rows stay out of sample.

\subsection{Moving block bootstrap}\label{en:sec:bootstrap}

With $B$ = \code{bootstrap\_samples} $>0$, the engine estimates the sampling distribution of
the selected model's coefficients. Each sample $b=1,\dots,B$:
\begin{enumerate}
  \item draws block starts $s$ uniformly in $\{0,\dots,n-\ell\}$, with $\ell$ =
        \code{bootstrap\_block\_size} $\le n$, and concatenates the blocks of rows
        $s,\dots,s+\ell-1$ until there are $n$ rows (the last block is truncated);
  \item refits the elastic net on those rows, for the selected factors, at the selected
        $(\alpha,\lambda)$;
  \item records the coefficients, the intercept and the MSE of that model on the test rows.
\end{enumerate}
With $\ell=1$ this is the ordinary i.i.d.\ bootstrap. Daily returns are not independent:
their volatility clusters. Blocks keep the dependence within $\ell$ consecutive rows, and the
moving block bootstrap is consistent for weakly dependent series when
$\ell\to\infty$ and $\ell/n\to0$ (K\"unsch, 1989). A block of about a month (20 rows) is a
reasonable choice for daily data.

\paragraph{Summaries.} For each selected factor, with draws $c_1,\dots,c_B$:
\begin{itemize}
  \item the mean, and the standard error $\sqrt{\frac1{B-1}\sum_b(c_b-\bar c)^2}$;
  \item the percentile interval at level $\gamma$ (\code{confidence\_level}): the empirical
        quantiles at $q=(1-\gamma)/2$ and $1-q$. With the sorted draws $c_{(0)}\le\dots\le
        c_{(B-1)}$, the quantile at $q$ interpolates linearly at position $q(B-1)$;
  \item the \emph{selection frequency}, the fraction of samples with $c_b\ne0$. It is
        meaningful for $\alpha>0$; for ridge every coefficient is non-zero, so it is always
        1;
  \item the intercept's mean and standard error, and the mean and percentile interval of
        the test MSE.
\end{itemize}

\paragraph{What the bootstrap does not cover.} The resamples refit the \emph{selected} subset
at the \emph{selected} $(\alpha,\lambda)$. The intervals are conditional on that choice: they
measure estimation noise, not the uncertainty of selecting the factors. Redoing the whole
search on every resample would cost $B$ times the search.

\paragraph{Reproducibility.} Sample $b$ uses its own Mersenne Twister, seeded with
\code{std::seed\_seq} from the 64-bit \code{bootstrap\_seed} and $b$. The draws do not depend
on which thread runs which sample.

\subsection{Parallelism and determinism}\label{en:sec:parallel}
Subset evaluations and bootstrap samples run on \code{nabia::detail::parallel\_for} with
\code{num\_threads} threads (0 = all cores). Every task writes only its own slot. Selections
use a sequential argmin with first-index tie-breaking, and random draws are seeded per
sample. Results are therefore bit-identical for every thread count, and a test checks this.

\subsection{Complexity}
One elastic net fit along an $L$-point path costs about $O(\text{sweeps}\cdot n|S|)$. One
subset evaluation costs $|\text{alphas}|\times(K+1)$ path fits. The search multiplies that by
the number of subsets, and the bootstrap adds $B$ single fits. In the factor replication
example (14 factors, sizes 1--4 exhaustive, 3 alphas $\times$ 30 lambdas, 5 folds, about
2\,600 training rows, 500 bootstrap samples), the whole example script, with its other
variants, runs in about 4\,s on a laptop.

%==============================================================================
\section{Implementation}

\subsection{Code map}
Paths are relative to the C++ project, \code{cpp/}.
\begin{center}\small
\begin{tabularx}{\linewidth}{@{}l X@{}}
\toprule
Concept & Code \\
\midrule
Parameters, fit, class & \nolinkurl{include/nabia/regression/ElasticNet.h}: \code{ElasticNetParams}, \code{ElasticNetFit}, \code{ElasticNet} \\
Centring, scaling (\S\ref{en:sec:standardize}) & \nolinkurl{src/regression/ElasticNet.cpp}: \code{standardize\_data} \\
Coordinate descent (Prop.~\ref{en:prop:cd}) & \code{CoordinateDescent::solve}, \code{CoordinateDescent::sweep}, \code{soft\_threshold} \\
$\lambda_{\max}$, grid (Prop.~\ref{en:prop:lmax}) & \code{nabia::regression::lambda\_max}, \code{lambda\_grid} \\
Path and warm starts & \code{nabia::regression::elastic\_net\_path} \\
Degrees of freedom (Prop.~\ref{en:prop:df}) & \code{degrees\_of\_freedom} \\
Back to original units & \code{make\_fit} \\
Engine parameters and results & \nolinkurl{include/nabia/regression/RegressionEngine.h} \\
Split, search, final model & \nolinkurl{src/regression/RegressionEngine.cpp}: \code{run\_regression\_engine} \\
Folds (\S\ref{en:sec:cv}) & \code{make\_folds} \\
CV, 1-SE rule, scores & \code{evaluate}, \code{in\_sample\_criterion}, \code{r\_squared} \\
Subset search (\S\ref{en:sec:search}) & \code{exhaustive\_search}, \code{forward\_search}, \code{count\_subsets} \\
Bootstrap (\S\ref{en:sec:bootstrap}) & \code{bootstrap}, \code{percentile} \\
Python bindings & \nolinkurl{bindings/src/bind_regression.cpp} \\
pandas interface & \nolinkurl{bindings/nabia_python_bridge/frames.py}: \code{regress} \\
Tests & \nolinkurl{tests/test_regression.cpp}, \nolinkurl{../nabia_python/tests/bridge/test_regression.py} \\
\bottomrule
\end{tabularx}
\end{center}

\subsection{Parameters}
\begin{center}\small
\begin{tabularx}{\linewidth}{@{}l l X@{}}
\toprule
Field (\code{RegressionEngineParams}) & Default & Meaning \\
\midrule
\code{min\_factors}, \code{max\_factors} & 1, 0 & subset sizes $k_{\min}..k_{\max}$; $k_{\max}=0$ means all \\
\code{search} & \code{Auto} & \code{Exhaustive}, \code{Forward} or \code{Auto} \\
\code{max\_exhaustive\_subsets} & 10\,000 & Auto: largest exhaustive search \\
\code{train\_fraction} & 0.7 & $f$; 1 = in sample \\
\code{in\_sample\_criterion} & \code{BIC} & \code{BIC}, \code{AIC} or \code{MSE} (when $f=1$) \\
\code{refit\_on\_full\_sample} & false & report coefficients refitted on all rows \\
\code{alphas} & 0, .25, .5, .75, 1 & mixes $\alpha$ to try \\
\code{lambdas} & empty & fixed grid; empty = automatic \\
\code{n\_lambdas}, \code{lambda\_min\_ratio} & 50, $10^{-3}$ & automatic grid \\
\code{validation} & \code{BlockedKFold} & or \code{WalkForward} \\
\code{validation\_folds}, \code{validation\_gap} & 5, 0 & $K$ and $g$ \\
\code{one\_standard\_error\_rule} & false & sparser model within 1 SE \\
\code{bootstrap\_samples}, \code{bootstrap\_block\_size} & 0, 1 & $B$ (0 = off) and $\ell$ \\
\code{confidence\_level}, \code{bootstrap\_seed} & 0.95, 42 & intervals and seed \\
\code{fit\_intercept}, \code{standardize} & true, true & see \S\ref{en:sec:standardize} \\
\code{tolerance}, \code{max\_iter} & $10^{-7}$, $10^5$ & solver stopping rule and sweep cap \\
\code{num\_threads} & 0 & 0 = all cores \\
\bottomrule
\end{tabularx}
\end{center}
Invalid values throw \code{std::invalid\_argument} at construction, with a message naming the
field. Python raises \code{ValueError}.

\subsection{Usage from C++}
\begin{lstlisting}
#include "nabia/regression/RegressionEngine.h"

auto params = std::make_unique<RegressionEngineParams>();
params->min_factors = 1;
params->max_factors = 4;
params->search = SubsetSearch::Exhaustive;
params->train_fraction = 0.7;                // last 30% of the rows out of sample
params->alphas = {0.0, 0.5, 1.0};
params->n_lambdas = 30;
params->validation = ValidationScheme::BlockedKFold;
params->bootstrap_samples = 500;
params->bootstrap_block_size = 20;           // keep about a month of dependence

RegressionEngine engine(std::move(params));
// Index-free: y is T, X is T x N (rows = time), names optional.
RegressionComputation r = engine.run(y, X, names);
r.factor_names; r.coefficients; r.alpha; r.lambda;  // the selected model
r.best.test_r2; r.bootstrap->lower; r.bootstrap->upper;

// With an index: target and factors must share it (align with drop_missing first).
RegressionResult<Date> rr = engine.run(target_series, factor_frame);
rr.predictions; rr.residuals; rr.test_start;

// A single elastic net, or a whole path.
ElasticNet en(std::make_unique<ElasticNetParams>(ElasticNetParams{.lambda = 1e-4, .alpha = 0.5}));
en.fit(X, y);
const ElasticNetFit& fit = en.result();      // coefficients in original units
auto grid = nabia::regression::lambda_grid(nabia::regression::lambda_max(X, y, 1.0), 60, 1e-4);
auto path = nabia::regression::elastic_net_path(X, y, ElasticNetParams{.alpha = 1.0}, grid);
\end{lstlisting}

\subsection{Usage from Python}
The module \code{nabia\_python\_bridge.regression} exposes the same classes, with
\code{lambda} spelled \code{lambda\_}. Arrays are \code{float64} numpy arrays, and the
computation runs with the GIL released. \code{frames.regress} adds pandas alignment and
labels:
\begin{lstlisting}[language=Python]
import numpy as np
from nabia_python_bridge import regression, frames

# numpy: y (T,), X (T, N)
r = regression.RegressionEngine(min_factors=1, max_factors=4, train_fraction=0.7,
                                search=regression.SubsetSearch.Exhaustive,
                                alphas=[0.0, 0.5, 1.0], n_lambdas=30,
                                bootstrap_samples=500, bootstrap_block_size=20).run(y, X, names)
r.factor_names, r.coefficients, r.alpha, r.lambda_
r.best.test_r2, r.validation_curve.mse, r.bootstrap.lower, r.bootstrap.upper

# In sample, chosen by BIC
regression.RegressionEngine(max_factors=5, train_fraction=1.0,
                            in_sample_criterion=regression.InSampleCriterion.BIC).run(y, X, names)

# pandas: aligned on the common index, rows with a missing value dropped
out = frames.regress(fund_returns, factor_returns, min_factors=1, max_factors=4,
                     validation="WalkForward", bootstrap_samples=500, bootstrap_block_size=20)
out.coefficients        # Series indexed by factor name
out.bootstrap           # DataFrame: mean, std_error, lower, upper, selection_frequency
out.best_by_size        # DataFrame: one row per subset size
out.predictions, out.residuals, out.test_start

# The elastic net itself
lmax = regression.lambda_max(X, y, alpha=1.0)
path = regression.elastic_net_path(X, y, regression.ElasticNetParams(alpha=1.0),
                                   regression.lambda_grid(lmax, 60, 1e-4))
\end{lstlisting}

\subsection{Worked example: recovering a known factor mix}
A former example of the project (factor replication, since removed) built a ``fund'' from real daily returns
(2012--2026):
\[
  \text{fund}_t=0.0001+0.50\,\text{NDX}_t+0.30\,\text{Gold}_t+0.20\,\text{10y}_t+\varepsilon_t,
  \qquad\varepsilon_t\sim\N\big(0,(0.02/\sqrt{252})^2\big).
\]
It then lets the engine choose 1 to 4 factors out of 14 candidates. The candidates include
close substitutes: the S\&P 500 ($\rho=0.93$ with the Nasdaq), silver, and 5y/30y bonds.
With a 70/30 split, 3 alphas, 30 lambdas, 5 blocked folds and 500 bootstrap samples in
20-day blocks:
\begin{center}\small
\begin{tabular}{@{}lccc@{}}
\toprule
Factor & True & Estimate & 95\% bootstrap interval \\
\midrule
Nasdaq 100 & 0.50 & 0.497 & [0.493, 0.500] \\
Gold & 0.30 & 0.300 & [0.295, 0.305] \\
10y notes & 0.20 & 0.193 & [0.175, 0.212] \\
Copper (not in the fund) & 0 & 0.003 & [$-0.002$, 0.007] \\
\bottomrule
\end{tabular}
\end{center}
\begin{itemize}
  \item The test $R^2$ is 0.981, and the out-of-sample tracking error is 1.98\%/yr against
        the 2\% noise.
  \item Copper wins the fourth slot on the test set by noise (the remark of
        \S\ref{en:sec:search}), and its interval contains zero.
\end{itemize}
The script plots the lasso path, the best $R^2$ by subset size, the validation curve with
its standard errors, the bootstrap intervals against the true weights, and the cumulative
fund against its replication.

%==============================================================================
\section{Verification}
\code{tests/test\_regression.cpp} (Catch2) checks:
\begin{itemize}
  \item $\lambda=0$ reproduces OLS;
  \item $\alpha=0$ reproduces the ridge closed form $(\tilde X^\T\tilde X+n\lambda I)^{-1}\tilde X^\T\tilde y$;
  \item solutions satisfy the optimality conditions \eqref{en:eq:kkt};
  \item $\lambda_{\max}$ zeroes every coefficient exactly, the next point of the grid does
        not, and the grid runs from $\lambda_{\max}$ to $r\lambda_{\max}$;
  \item constant factors get a zero coefficient;
  \item exhaustive and forward search recover the true factors out of sample, and Auto
        switches to Forward above the limit;
  \item in-sample BIC rejects noise factors;
  \item a single-point grid skips validation;
  \item walk-forward validation works with the 1-SE rule;
  \item bootstrap intervals cover the true coefficients;
  \item results, bootstrap draws included, are identical for 1 and 4 threads;
  \item refitting keeps the out-of-sample predictions;
  \item the \code{DataFrame} interface keeps names and the index;
  \item invalid parameters and data throw.
\end{itemize}
\code{nabia\_python/tests/bridge/test\_regression.py} repeats the key checks through Python.
The suites also run under AddressSanitizer.

\section{Limitations}
\begin{itemize}
  \item The model is linear with constant coefficients. Time-varying exposures need rolling
        windows, or a state-space model such as the Kalman filter of the forecasting module.
  \item The bootstrap is conditional on the selected subset and hyperparameters
        (\S\ref{en:sec:bootstrap}).
  \item Exhaustive search grows as $\binom Nk$. Forward search is greedy.
  \item The test error of the selected model is optimistic when many subsets are compared.
\end{itemize}

\section*{References}
\begin{itemize}\small
  \item J.~Friedman, T.~Hastie, R.~Tibshirani (2010). Regularization paths for generalized linear
        models via coordinate descent. \emph{Journal of Statistical Software} 33(1).
  \item H.~Zou, T.~Hastie (2005). Regularization and variable selection via the elastic net.
        \emph{JRSS B} 67(2).
  \item H.~Zou, T.~Hastie, R.~Tibshirani (2007). On the ``degrees of freedom'' of the lasso.
        \emph{Annals of Statistics} 35(5).
  \item P.~Tseng (2001). Convergence of a block coordinate descent method for nondifferentiable
        minimization. \emph{JOTA} 109(3).
  \item B.~Efron (2004). The estimation of prediction error: covariance penalties and
        cross-validation. \emph{JASA} 99(467).
  \item H.~K\"unsch (1989). The jackknife and the bootstrap for general stationary observations.
        \emph{Annals of Statistics} 17(3).
  \item J.~Campbell, S.~Thompson (2008). Predicting excess stock returns out of sample.
        \emph{Review of Financial Studies} 21(4).
  \item T.~Hastie, R.~Tibshirani, J.~Friedman (2009). \emph{The Elements of Statistical Learning},
        2nd ed., Springer.
\end{itemize}
